% Part: propositional-logic
% Chapter: propositional-logic
% Section: soundness

\documentclass[../../../include/open-logic-section]{subfiles}

\begin{document}

\olfileid{pl}{prp}{snd}

\olsection{বচনীয় যুক্তিবিদ্যার বিশুদ্ধতা}

\begin{thm}[বিশুদ্ধতা]
\ollabel{thm:soundness}
$\Gamma\Proves!A$ হলে
$\Gamma \Entails !A$।
\end{thm}

\begin{proof}
পূর্বধারণাসমেত প্রমাণের গঠনের উপর আরোহ করি।
$!A$ স্বতঃসিদ্ধ হলে $\Entails !A$; সুতরাং $\Gamma \Entails!A$-ও হয়।
একইভাবে, $!A \in \Gamma$ হলে $\Gamma \Entails!A$।
আরোহী ধাপে ধরা যাক $!C$ ও $!C\lif!A$ থেকে
\emph{মোডাস পোনেন্স} দ্বারা $!A$ পাওয়া যায়। আরোহের অনুমান অনুসারে
$!C\lif !A$ ও $!C$-র জন্য উপপাদ্যটি সিদ্ধ।
\olref[pl][syn][sem]{prop:semanticalfacts}-এর দ্বিতীয় দফা থেকে
প্রয়োজনীয় ফল পাওয়া যায়।
\end{proof}

\begin{cor}
$\Gamma$ পরিতৃপ্তিযোগ্য হলে $\Gamma$ সঙ্গত।
তাই বচনীয় যুক্তিবিদ্যা সঙ্গত।
\end{cor}
% BN-SRC-676 TARGET-CORRECTION-BEGIN
\par\smallskip\noindent\textnormal{\small\textbf{উৎস-সংশোধন (BN-SRC-676):} উপপাদ্যটি পূর্বধারণা থেকে প্রমাণের জন্য; তাই উৎসের “উপপাদ্যের উপর আরোহ” কথাটি প্রমাণের গঠনের উপর আরোহ হিসেবে স্পষ্ট করা হয়েছে।} % BN-SRC-676 TARGET-CORRECTION-END
% BN-SRC-677 TARGET-CORRECTION-BEGIN
\par\smallskip\noindent\textnormal{\small\textbf{উৎস-সংশোধন (BN-SRC-677):} অর্থগত মোডাস পোনেন্স সংশ্লিষ্ট প্রস্তাবের দ্বিতীয় দফা, তৃতীয়টি নয়। নির্দেশে সঠিক অংশ/অধ্যায়/বিভাগও বসানো হয়েছে।} % BN-SRC-677 TARGET-CORRECTION-END

\end{document}
